\section{Irreducible spaces. Noetherian spaces}
\label{section:0.2}

\subsection{Irreducible spaces}
\label{subsection:0.2.1}

\begin{env}[2.1.1]
\label{0.2.1.1}
A topological space $X$ is said to be \emph{irreducible} if it is non-empty
and is not the union of two closed subspaces distinct from $X$. It
amounts to the same to say that $X\neq\varnothing$ and that the intersection of two
(and hence of finitely many) non-empty open subsets of $X$ is non-empty,
or that every non-empty open subset is everywhere dense, or that every closed
subset $\neq X$ is \emph{nowhere dense}, or finally that every open subset of $X$ is
\emph{connected}.
\end{env}

\begin{env}[2.1.2]
\label{0.2.1.2}
For a subspace $Y$ of a topological space $X$ to be irreducible, it is
necessary and sufficient that its closure $\overline Y$ be irreducible. In
particular, every subspace which is the closure $\overline{\{x\}}$ of a
subspace reduced to one point is irreducible; we shall express the relation
$y\in\overline{\{x\}}$ (equivalent to
$\overline{\{y\}}\subset\overline{\{x\}}$) by saying that $y$ is a
\emph{specialization of $x$} or that $x$ is a \emph{generization of $y$}.
When there exists in an irreducible space $X$ a point $x$ such that
$X=\overline{\{x\}}$, we shall say that $x$ is a \emph{generic point} of $X$.
Every non-empty open subset of $X$ then contains $x$, and every subspace containing
$x$ has $x$ as a generic point.
\end{env}

\begin{env}[2.1.3]
\label{0.2.1.3}
Recall that a \emph{Kolmogorov space} is a topological space $X$
satisfying the separation axiom:

\smallskip
\noindent
\textup{(T$_0$)} If $x\neq y$ are any two points of $X$, there exists
an open set containing one of the points $x$, $y$ and not the other.

If an irreducible Kolmogorov space has a generic point, it has
\emph{only one}, since a non-empty open set contains every generic
point.

Recall that a topological space $X$ is said to be \emph{quasi-compact} if from every
open covering of $X$ one can extract a finite covering of $X$ (or,
what amounts to the same, if every decreasing filtered family of non-empty closed
sets has a non-empty intersection). If $X$ is a
quasi-compact space, every non-empty closed subset $A$ of $X$ contains a
\emph{minimal} non-empty closed set $M$, since the set of non-empty closed subsets
of $A$ is inductive for the relation $\supset$; if moreover $X$ is a
Kolmogorov space, $M$ is necessarily reduced to a single point (or, as one
says by abuse of language, is a \emph{closed point}).
\end{env}

\begin{env}[2.1.4]
\label{0.2.1.4}
In an irreducible space $X$, every non-empty open subspace $U$ is
irreducible, and if $X$ has a generic point $x$, $x$ is also a generic
point of $U$.

Let $(U_\alpha)$ be a covering (with non-empty index set) of a
topological space $X$, consisting of non-empty open sets; for $X$ to be
irreducible, it is necessary and sufficient that $U_\alpha$ be irreducible for every
$\alpha$, and that $U_\alpha\cap U_\beta\neq\varnothing$ for all
$\alpha$, $\beta$. The condition is obviously necessary; to see that it
is sufficient, it suffices to prove that if $V$ is a non-empty open subset of $X$,
$V\cap U_\alpha$ is non-empty for every $\alpha$, for then $V\cap U_\alpha$
is dense in $U_\alpha$ for every $\alpha$, and consequently $V$ is dense in
$X$. Now, there is at least one index $\gamma$ such that
$V\cap U_\gamma\neq\varnothing$, so $V\cap U_\gamma$ is dense in
$U_\gamma$, and since, for every $\alpha$,
$U_\alpha\cap U_\gamma\neq\varnothing$, we also have
$V\cap U_\alpha\cap U_\gamma\neq\varnothing$.
\end{env}

\oldpage[0\textsubscript{I}]{22}

\begin{env}[2.1.5]
\label{0.2.1.5}
Let $X$ be an irreducible space and $f$ a continuous map from $X$ to a
topological space $Y$. Then $f(X)$ is irreducible, and if $x$ is a
generic point of $X$, $f(x)$ is a generic point of $f(X)$ and consequently also of
$\overline{f(X)}$. In particular, if moreover $Y$ is irreducible and has only one
generic point $y$, then for $f(X)$ to be everywhere dense it is necessary and sufficient
that $f(x)=y$.
\end{env}

\begin{env}[2.1.6]
\label{0.2.1.6}
Every irreducible subspace of a topological space $X$ is contained in a
\emph{maximal} irreducible subspace, which is necessarily closed. The
maximal irreducible subspaces of $X$ are called the \emph{irreducible
components} of $X$. If $Z_1$, $Z_2$ are two distinct irreducible components
of the space $X$, $Z_1\cap Z_2$ is a \emph{nowhere dense} closed set
in each of the subspaces $Z_1$, $Z_2$; in particular, if an irreducible
component of $X$ has a generic point (2.1.2), such a point cannot
belong to any other irreducible component. If $X$ has only a
\emph{finite} number of irreducible components $Z_i$ ($1\leq i\leq n$), and if, for
each $i$, we set
\[
  U_i=Z_i\cap\complement\Bigl(\bigcup_{j\neq i}Z_j\Bigr),
\]
then the $U_i$ are open, irreducible and pairwise disjoint, and their
union is dense in $X$.

Let $U$ be an open subset of a topological space $X$. If $Z$ is an
irreducible subset of $X$ meeting $U$, $Z\cap U$ is open and dense in $Z$,
hence irreducible; conversely, for every irreducible closed subset $Y$ of
$U$, the closure $\overline Y$ of $Y$ in $X$ is irreducible and
$\overline Y\cap U=Y$. One concludes that there is a \emph{bijective
correspondence} between the irreducible components of $U$ and the irreducible
components of $X$ which meet $U$.
\end{env}

\begin{env}[2.1.7]
\label{0.2.1.7}
If a topological space $X$ is the union of a \emph{finite} number of
irreducible closed subspaces $Y_i$, the irreducible components of $X$
are the maximal elements of the set of the $Y_i$, for if $Z$ is an irreducible
closed subset of $X$, $Z$ is the union of the $Z\cap Y_i$, from which one deduces
at once that $Z$ must be contained in one of the $Y_i$. Let $Y$ be a subspace
of a topological space $X$, and suppose that $Y$ has only a finite number of
irreducible components $Y_i$ ($1\leq i\leq n$); then the closures
$\overline{Y_i}$ in $X$ are the irreducible components of $\overline Y$.
\end{env}

\begin{env}[2.1.8]
\label{0.2.1.8}
Let $Y$ be an irreducible space having only one generic point $y$. Let
$X$ be a topological space and $f$ a continuous map from $X$ to $Y$.
Then, for every irreducible component $Z$ of $X$ meeting $f^{-1}(y)$,
$f(Z)$ is dense in $Y$. The converse is not necessarily true;
however, if $Z$ has a generic point $z$ and $f(Z)$ is dense in
$Y$, we necessarily have $f(z)=y$ (2.1.5); moreover, $Z\cap f^{-1}(y)$ is
then the closure of $\{z\}$ in $f^{-1}(y)$ and is therefore irreducible, and
since every irreducible subset of $f^{-1}(y)$ containing $z$ is
necessarily contained in $Z$ (2.1.6), $z$ is a generic point of
$Z\cap f^{-1}(y)$. Since every irreducible component of $f^{-1}(y)$ is
contained in an irreducible component of $X$, we see that if every
irreducible component $Z$ of $X$ meeting $f^{-1}(y)$ has a generic
point, then there is a \emph{bijective correspondence} between the set of
these components and the set of irreducible components
$Z\cap f^{-1}(y)$ of $f^{-1}(y)$, the generic points of $Z$ being
identical with those of $Z\cap f^{-1}(y)$.
\end{env}

\oldpage[0\textsubscript{I}]{23}

\subsection{Noetherian spaces}
\label{subsection:0.2.2}

\begin{env}[2.2.1]
\label{0.2.2.1}
A topological space $X$ is said to be \emph{Noetherian} if the set of
open subsets of $X$ satisfies the \emph{maximal} condition, or, what amounts to the
same, if the set of closed subsets of $X$ satisfies the \emph{minimal} condition. We
say that $X$ is \emph{locally Noetherian} if every $x\in X$ has a
neighbourhood which is a Noetherian subspace.
\end{env}

\begin{env}[2.2.2]
\label{0.2.2.2}
Let $E$ be an ordered set satisfying the \emph{minimal} condition, and let
$P$ be a property of the elements of $E$ subject to the following condition: if
$a\in E$ is such that $P(x)$ is true for every $x<a$, then $P(a)$ is
true. Under these conditions, \emph{$P(x)$ is true for every $x\in E$}
(``principle of Noetherian induction''). Indeed, let $F$ be the set of
$x\in E$ for which $P(x)$ is false; if $F$ were non-empty, it would have a
minimal element $a$, and since $P(x)$ is then true for every $x<a$, $P(a)$
would also be true, which is a contradiction.

We shall apply this principle in particular when $E$ is a \emph{set
of closed subsets of a Noetherian space}.
\end{env}

\begin{env}[2.2.3]
\label{0.2.2.3}
Every subspace of a Noetherian space is Noetherian. Conversely, every
topological space which is a finite union of Noetherian subspaces is Noetherian.
\end{env}

\begin{env}[2.2.4]
\label{0.2.2.4}
Every Noetherian space is quasi-compact; conversely, every topological space
in which every open subset is quasi-compact is Noetherian.
\end{env}

\begin{env}[2.2.5]
\label{0.2.2.5}
A Noetherian space has only a \emph{finite} number of irreducible
components, as one sees by Noetherian induction.
\end{env}

